\documentclass[12pt,a4paper,oneside]{article}
\usepackage[brazil]{babel}
\usepackage[utf8]{inputenc}
\usepackage{olymp}
\usepackage[dvipsnames]{xcolor}

\setcounter{problem}{0}

\begin{document}

\begin{problem}{Maior dígito}{digito}{2}

Faça um programa para ler dois números inteiros de 4 dígitos e calcular sua soma, mas, em vez de escrever esses valores, escrever apenas o dígito de maior valor em cada um deles.

Por exemplo, para $2026$ e $1926$ temos $2026 + 1926 = 3952$, então o programa deve escrever $6$, $9$ e $9$,
respectivamente o dígito de maior valor de $2026$, $1926$ e $3952$.

A fim de evitar repetição de código, seu programa deve usar obrigatoriamente uma função com o seguinte protótipo:

\texttt{\textcolor{Mulberry}{int} maiordigito(\textcolor{Mulberry}{int} n)}

A função recebe como parâmetro um número inteiro de 4 dígitos e retorna o valor do maior dígito do número.

%\Notes
%
%Observações, se tiver.

\InputFile

A entrada contém dois números inteiros $A$ e $B$. Restrições: $1000 \leq A, B, A+B \leq 9999$.

\OutputFile

A saída deve conter 3 linhas com respectivamente o maior dígito de $A$, de $B$ e de $A+B$.

%\Notes

\Example

\begin{example}
\exmp{
2026 1926
}{
6
9
9
}%
\end{example}

\begin{example}
\exmp{
1234 5678
}{
4
8
9
}%
\end{example}

\begin{example}
\exmp{
1000 5000
}{
1
5
6
}%
\end{example}

%Comentários sobre o exemplo, se necessário.

\end{problem}

\end{document}